\chapter{Objetivos}
\label{chap:objetivos}

\section{Objetivo general}

El análisis de muestras de lenguaje (\emph{Language Sample Analysis}, \ac{LSA}) constituye el
patrón de referencia para la evaluación expresiva en intervención logopédica, y el
protocolo para el análisis de muestras de lenguaje asistido por \ac{SAAC}, definido por Gloria
Soto~\cite{soto2020protocolo}, traslada esa metodología al caso de los usuarios de
Sistemas de Comunicación Aumentativa y Alternativa (\ac{SAAC}). Su aplicación manual exige
segmentar y codificar al menos cincuenta enunciados a lo largo de cuatro dimensiones, 
lo que se traduce en un coste de varias horas de trabajo profesional por sesión grabada. 
Este coste temporal limita en la práctica la frecuencia con la que los especialistas pueden 
aplicar el protocolo, y motiva la búsqueda de un sistema que reduzca dicha carga sin sustituir 
el juicio clínico del profesional.

El objetivo general de este trabajo es diseñar e implementar un prototipo de
automatización parcial del protocolo \ac{PAALSS} centrado en el \textbf{canal auditivo}
de la sesión, mediante modelos de \ac{IA} y \ac{PLN} que reduzcan la carga de trabajo
manual del especialista. El sistema actúa como asistente de pre-anotación: propone
valores para las métricas del protocolo a partir del procesamiento del audio de la
sesión, mientras que la validación final recae siempre en el profesional, sin que el
sistema emita diagnóstico alguno ni sustituya el juicio clínico experto. El análisis
del canal visual así como del registro del dispositivo (selección de pictogramas y gestos) 
se plantea como extensión futura del prototipo, una vez validado el canal auditivo sobre 
las sesiones disponibles.

Ningún trabajo previo revisado en el estado del arte aborda esta automatización para el
protocolo \ac{PAALSS}. Los precedentes más cercanos operan sobre transcripciones ya elaboradas, 
en inglés y sin componente multimodal, por lo que no resuelve ni la extracción de enunciados
desde vídeo ni el pre-rellenado de un formulario clínico estructurado en español. Del
mismo modo, las herramientas de \ac{LSA} automático de mayor difusión (Batchalign, 
\ac{SALT}, \ac{CLAN}) se han desarrollado y validado sobre el inglés, y los modelos de \ac{PLN} 
disponibles para el español, como \emph{BETO} o \emph{spaCy-es}, degradan su rendimiento sobre 
habla espontánea e infantil. Esta ausencia de precedente directo justifica el planteamiento de los
cuatro objetivos específicos que se describen a continuación.

Por último, cabe destacar que el sistema se concibe bajo una restricción de procesamiento
local: dado que los datos manejados corresponden a sesiones de menores con discapacidad,
ninguna fase del \emph{pipeline} puede depender de servicios en la nube. 
Esta restricción condiciona de forma transversal el diseño de todos los objetivos específicos y, 
en particular, la elección de las tecnologías empleadas en cada módulo.

\section{Objetivos específicos}

A partir del objetivo general, se han definido cuatro objetivos específicos que
cubren de forma secuencial las fases necesarias para construir el sistema: extracción de
enunciados desde el canal auditivo del vídeo, procesamiento lingüístico de dichos
enunciados, clasificación morfológica de las categorías ambiguas propias del lenguaje
asistido, y generación automática del pre-rellenado del formulario. Cada uno de ellos
responde a una carencia concreta detectada en la revisión del estado del arte.

\subsection{OE1: Extracción automática de enunciados}

Implementar un sistema automático de extracción de enunciados desde el canal auditivo del
vídeo de la sesión, mediante transcripción automática del habla, diarización de hablantes
y fusión de ambas señales en Turnos Comunicativos (CT) atribuidos al hablante
correspondiente. La transcripción se apoya en modelos \emph{Whisper}~\cite{whisper_radford2023}
ejecutados de forma local mediante \texttt{faster-whisper}, mientras que la diarización
separa los turnos de la persona usuaria de SAAC de los de su interlocutor mediante
enfoques como \emph{pyannote.audio}~\cite{pyannote_plaquet2023} o
\emph{DiCoW}~\cite{dicow_diarization}.
El análisis de la modalidad visual, selección de pictogramas, gesto y mirada, queda planteado
como línea de trabajo futuro, una vez implementado la solución para el canal auditivo.

\subsection{OE2: Pipeline de procesamiento lingüístico}

Diseñar un \emph{pipeline} de procesamiento lingüístico que identifique \textbf{categorías
gramaticales}, \textbf{estructuras morfológicas} y \emph{patrones sintácticos} definidos en el 
protocolo de Gloria Soto, adaptando los modelos de \ac{PLN} disponibles para el español a las
particularidades del lenguaje producido con \ac{SAAC}. Estas particularidades incluyen el uso
de pictogramas en forma de cita (infinitivo, singular), la ausencia de inflexión
morfológica explícita en la selección y la naturaleza telegráfica de los enunciados. Dado
que herramientas como \emph{spaCy}~\cite{spoken_spanish_pos} o
\emph{BETO}~\cite{beto_morphological} se han entrenado sobre habla y texto estándar,
su aplicación directa requiere una capa de adaptación que tenga en cuenta el contexto
multimodal de la sesión y no únicamente la etiqueta léxica del pictograma.

\subsection{OE3: Modelos de análisis gramatical}

Desarrollar un clasificador de las categorías morfológicas ambiguas propias del
lenguaje asistido por pictogramas: aquellas que no pueden resolverse mediante el
etiquetado estándar de herramientas de \ac{PLN} entrenadas sobre habla oral, porque
el contexto comunicativo del \ac{SAAC} altera el valor morfológico esperado de la
forma de cita. La estrategia adoptada combina reglas lingüísticas explícitas para
los casos deterministas con clasificación \emph{zero-shot} mediante el modelo de
lenguaje \emph{BETO}~\cite{beto_morphological} y \emph{Pattern-Exploiting
Training}~\cite{fewshot_clinical_nlp} para los casos en que el contexto es
necesario para dirimir la ambigüedad.

\subsection{OE4: Módulo de pre-rellenado del protocolo}

Desarrollar un módulo que traduzca las métricas calculadas en OE2 y OE3 en propuestas
de valores para cada campo del formulario \ac{PAALSS}, de forma trazable y vinculada
a los enunciados que las sustentan. Sumado a esto, se propone una arquitectura de generación
aumentada por recuperación (\emph{Retrieval-Augmented Generation}, \ac{RAG})
\cite{rag_lewis2020} que recupera los \ac{CT} semánticamente relevantes para cada
campo mediante una base vectorial, y los emplea como contexto para un modelo de
lenguaje local que genera observaciones acerca de la sesión. Este \emph{RAG} no representa una 
arquitectura completa final, sino que se plantea como una primera versión parcial para un desarrollo
futuro. El resultado se exporta como \ac{PDF} editable del formulario oficial, de modo que el especialista pueda revisar,
corregir y validar cada campo propuesto sin necesidad de partir de cero.

\section{Alcance y limitaciones}

El presente trabajo se circunscribe al español y a los tableros de comunicación basados en
el catálogo de pictogramas de \textbf{\ac{ARASAAC}}, y se valida únicamente sobre dos sesiones
anotadas por expertos correspondientes a una misma usuaria, lo que condiciona la
generalización de los resultados a otros perfiles de usuario o a otros sistemas de
comunicación aumentativa. El sistema se limita, además, a la generación de un
pre-rellenado del protocolo y no constituye en ningún caso una herramienta de diagnóstico,
quedando la validación final del formulario siempre a cargo del profesional. El
procesamiento de la modalidad visual (selección de pictogramas, gesto y mirada) se
plantea como línea de trabajo futuro y queda fuera del alcance de esta primera versión, que
se centra en el canal auditivo de la sesión. Por último, y dado que los datos manejados
corresponden a menores con discapacidad, todo el procesamiento se ejecuta de forma local,
sin recurrir a servicios en la nube, lo que excluye del alcance del trabajo cualquier
modelo o API que requiera el envío de datos a terceros.

\section{Requisitos}

En cuanto a los requisitos \emph{funcionales}, se pueden enumerar los siguientes:

\begin{itemize}
  \item \emph{Extracción de enunciados}. El sistema debe generar, a partir de un vídeo de
  sesión, una lista de \textbf{Turnos Comunicativos} con su transcripción, su atribución de
  hablante y sus marcas temporales.
  \item \emph{Procesamiento lingüístico}. Cada \textbf{Turno Comunicativo} debe poder anotarse con
  sus categorías gramaticales, su estructura morfológica y su patrón sintáctico.
  \item \emph{Pre-rellenado del protocolo}. El sistema debe producir, para cada sesión
  procesada, un conjunto de métricas y tablas correspondientes a las cuatro dimensiones del
  protocolo \textbf{\ac{PAALSS}}, trazables hasta los enunciados que las originan.
  \item \emph{Exportación}. El resultado del pre-rellenado debe poder exportarse en un
  formato que permita su revisión y edición por parte del profesional.
\end{itemize}

Por otra parte, los requisitos \emph{no funcionales} considerados son los siguientes:

\begin{itemize}
  \item \emph{Procesamiento local}. Ninguna fase del sistema puede depender de servicios
  en la nube, en cumplimiento del Reglamento General de Protección de Datos (RGPD) y dado
  que los datos corresponden a menores con discapacidad.
  \item \emph{Trazabilidad}. Toda métrica generada debe poder vincularse al enunciado o
  conjunto de enunciados que la sustentan, de forma que el profesional pueda auditar el
  resultado propuesto por el sistema.
  \item \emph{Idioma}. El sistema debe operar sobre transcripciones y enunciados en
  español, incluyendo las variantes propias del habla infantil y del lenguaje asistido.
  \item \emph{Validez clínica}. El sistema no sustituye el juicio profesional; sus salidas
  se conciben como propuesta de pre-anotación sujeta a confirmación o corrección por parte
  del especialista. No hay validación con ninguna métrica, debido a la falta de muestras y
  que no entra en los objetivos de este análisis.
\end{itemize}

% Local Variables:
%  coding: utf-8
%  mode: latex
%  mode: flyspell
%  ispell-local-dictionary: "castellano8"
% End: